Self-critique prompting is widely used to improve large language model reasoning, but it is unclear whether it produces meaningful internal state updates or only a second sample from the same policy. We study this question by jointly measuring behavior and residual-stream representation drift under a controlled evaluation protocol. On 48 examples (24 GSM8K and 24 CommonsenseQA), we compare four conditions with \qwen: one-shot baseline, re-decode control, intrinsic self-critique revision, and an external-critique subset using \gptfourone. We extract per-layer final-token residual representations and quantify drift with cosine distance and L2 displacement, then test condition differences with paired non-parametric statistics and FDR correction.

Self-critique substantially reduced accuracy from 39.6\% to 16.7\% (delta $-22.9$ points, 95\% CI $[-37.5,-6.25]$, McNemar $p=0.0098$). Re-decoding also reduced accuracy to 25.0\%, indicating broad second-pass instability. Critically, mean cosine drift was nearly identical for self-critique and re-decode (0.5700 vs 0.5719), and layer-wise differences were non-significant after Benjamini--Hochberg correction across 28 layers. External critique did not recover performance on its subset.

These results provide no evidence, in this setup, for uniquely structured reflection-induced residual drift. Instead, self-reflection behaves closer to destabilizing re-generation than targeted internal correction.
